\chapter{a} % Capítulo
\textbf{Esse template possui as configurações necessárias para um TCC/Monografia segundo a FAETERJ-Petrópolis, abaixo encontra-se um apanhado de funções que podem ser utilizadas na construção do documento.}


--------------------


Para referenciar algo basta utilizar \cite{unesco} %unesco foi o nome dado a uma referencia no arquivo monografia.tex

\section{Subseção} % Seção
Toda frase em inglês precisa ser italica \textit{Game}.



\begin{quote}
\small "Para citar uma frase de algum autor basta utilizar isso." 

(Autor)
\end{quote}



Para referenciar alguma figura como a logomarca da FAETERJ basta utilizar \ref{fig:faeterj}. % deve-se referenciar imagem antes que a mesma venha a ser colocada

\begin{figure}[!htb]
     \centering
     \includegraphics[scale=0.24]{imagens/faeterj.png}
     \caption{\small Logomarca da FAETERJ.}
     \label{fig:faeterj}
\end{figure}



Para criar uma lista de elementos basta utilizar o que está abaixo

\begin{itemize}
\item item1
\item item2
\item item3
\end{itemize}



Para criar uma equação basta utilizar \ref{eq:delta} ou $\fontsize{14}{14}{A}_{1}$.

\begin{equation}
	\fontsize{14}{14}{A}_{1} = {A}_{2} - {A}_{3}
	\label{eq:delta}
\end{equation}

\begin{equation}
	\fontsize{14}{14}{A}_{1} = \frac{{\alpha}^{\beta}}{\omega}
	\label{eq:delta}
\end{equation}



Para criar um algorimo basta utilizar \ref{alg:a}.

\begin{algorithm}[H]
float A;\\

\While{true}{	
	A = 0;\\

	\eIf{A > 60}{
		A = 30;\\
	}{
		A = 20;\\
   	}
}
\caption{Algoritmo criado.}
\label{alg:a}
\end{algorithm}



Para criar uma tabela basta utilizar \ref{tab:1}

\begin{table}[!htbp]
\centering
\caption{Tabela 1}
\label{tab:1}
\begin{tabular}{rlr}
 
Hora & Localidade & Resultado \\
\hline                             
12:00			&	Brazil/MG	&	True	\\
13:00			&	Brazil/SP	&	False	\\
14:00			&	Brazil/MG	&	False	\\
15:00			&	Brazil/RJ	&	True	\\
\end{tabular}
\end{table}

\textbf{Editado por:} Johnathan Fercher da Rosa
